Input to the Node Property Flow (NPF) Package is read from the file that has type ``NPF6'' in the Name File.  A single NPF Package is required for each GWF model. 

\vspace{5mm}
\subsubsection{Structure of Blocks}
\lstinputlisting[style=blockdefinition]{./mf6ivar/tex/gwf-npf-options.dat}
\lstinputlisting[style=blockdefinition]{./mf6ivar/tex/gwf-npf-griddata.dat}

\vspace{5mm}
\subsubsection{Explanation of Variables}
\begin{description}
\input{./mf6ivar/tex/gwf-npf-desc.tex}
\end{description}

\vspace{5mm}
\subsubsection{Comparison with MODFLOW-NWT}
The default harmonic-mean method and the LOGARITHMIC and AMT-LMK options of ALTERNATIVE\_CELL\_AVERAGING correspond to LAYAVG values of 0, 1, and 2 in the MODFLOW-NWT Upstream Weighting (UPW) Package \citep{modflownwt}. They calculate the same horizontal conductance as MODFLOW-NWT for confined cells and, when the Newton-Raphson formulation is used, for a convertible cell connected to a cell of the same thickness, apart from differences in how the saturation is smoothed. Where the thickness of connected cells differs, the conductance of a convertible cell differs, because MODFLOW-NWT calculates it from the saturated thickness of the upstream cell, while \mf multiplies a conductance calculated from the full thickness of both cells by the saturated fraction of the upstream cell.

\vspace{5mm}
\subsubsection{Example Input File}
\lstinputlisting[style=inputfile]{./mf6ivar/examples/gwf-npf-example.dat}

